\setchoices{label=(\Alph*)}

\begin{question}[index=1]
设 $z=f(x,y)$ 具有二阶连续偏导数，$x=u$，$y^2=v-u^2$，以 $u,v$ 为新的自变量变换方程
$\displaystyle y\frac{\partial z}{\partial x}-x\frac{\partial z}{\partial y}=0\ (y\ne0)$，得到
\begin{choices}[columns=4,label-pos=left]
\item $\dfrac{\partial^2 z}{\partial v^2}\ne0$．
\item $\dfrac{\partial^2 z}{\partial v\partial u}\ne0$．
\item $\dfrac{\partial z}{\partial v}=0$．
\item $\dfrac{\partial z}{\partial u}=0$．
\end{choices}
\end{question}
\QuestionBreak

\begin{question}[index=2]
已知 $f(x)=\mathrm e^{|x|}$，则 $f(x)$ 的原函数为
\begin{choices}[columns=2,label-pos=left]
\item $F(x)=\begin{cases}-\mathrm e^{-x}+C+2,&x<0,\\ \mathrm e^x+C,&x\ge0.\end{cases}$
\item $F(x)=\begin{cases}\mathrm e^{-x}+C,&x<0,\\ \mathrm e^x+C+2,&x\ge0.\end{cases}$
\item $F(x)=\begin{cases}-\mathrm e^{-x}+C,&x<0,\\ \mathrm e^x+C+2,&x\ge0.\end{cases}$
\item $F(x)=\begin{cases}-\mathrm e^x+C+2,&x<0,\\ \mathrm e^x+C,&x\ge0.\end{cases}$
\end{choices}
\end{question}
\QuestionBreak

\begin{question}[index=3]
设 $\displaystyle f(x)=\begin{cases}x,&0\le x\le\dfrac12,\\[3pt]2-2x,&\dfrac12<x<1,\end{cases}$
$\displaystyle S(x)=\frac{a_0}{2}+\sum_{n=1}^{\infty}a_n\cos n\pi x$，$-\infty<x<+\infty$，
其中 $\displaystyle a_n=2\int_0^1 f(x)\cos n\pi x\dd x$，则 $\displaystyle S\left(-\frac52\right)=$
\begin{choices}[columns=4,label-pos=left]
\item $\dfrac12$．
\item $-\dfrac12$．
\item $\dfrac34$．
\item $-\dfrac34$．
\end{choices}
\end{question}
\QuestionBreak

\begin{question}[index=4]
已知 $F(x,y,z)=f(\pi y-\sqrt2z)-\mathrm e^{2x-z}$，其中 $f$ 可微。
若 $aF'_x+bF'_y+cF'_z=0$，则 $(a,b,c)$ 可以是
\begin{choices}[columns=2,label-pos=left]
\item $\left(2,\dfrac{2\sqrt2}{\pi},1\right)$．
\item $\left(1,\dfrac{2\sqrt2}{\pi},2\right)$．
\item $\left(1,2,\dfrac{2\sqrt2}{\pi}\right)$．
\item $\left(2,1,\dfrac{2\sqrt2}{\pi}\right)$．
\end{choices}
\end{question}
\QuestionBreak

\begin{question}[index=5]
对 $3$ 阶矩阵 $\symbfit A$ 的伴随矩阵 $\symbfit A^*$ 先交换第 $1$ 行和第 $3$ 行，
再将第 $2$ 列的 $-2$ 倍加到第 $3$ 列，得到矩阵 $-\symbfit E$，$\symbfit E$ 为 $3$ 阶单位矩阵，则 $\symbfit A=$
\begin{choices}[columns=2,label-pos=left]
\item $\begin{pmatrix}0&0&1\\0&1&-2\\1&0&0\end{pmatrix}$ 或 $\begin{pmatrix}0&0&-1\\0&-1&2\\-1&0&0\end{pmatrix}$．
\item $\begin{pmatrix}0&0&1\\0&1&-2\\1&0&0\end{pmatrix}$ 或 $\begin{pmatrix}0&0&1\\-2&1&0\\1&0&0\end{pmatrix}$．
\item $\begin{pmatrix}0&0&-1\\0&-1&2\\-1&0&0\end{pmatrix}$ 或 $\begin{pmatrix}0&0&-1\\2&-1&0\\-1&0&0\end{pmatrix}$．
\item $\begin{pmatrix}0&0&1\\-2&1&0\\1&0&0\end{pmatrix}$ 或 $\begin{pmatrix}0&0&-1\\2&-1&0\\-1&0&0\end{pmatrix}$．
\end{choices}
\end{question}
\QuestionBreak

\begin{question}[index=6]
设 $3$ 阶实对称矩阵 $\symbfit A$ 的各行元素之和均为 $2$，其主对角线元素之和为 $5$，
$r(\symbfit A)=2$，则二次型 $f(x_1,x_2,x_3)=\symbfit x^{\mathrm T}\symbfit A\symbfit x$
在条件 $x_1^2+x_2^2+x_3^2=1$ 下的最大值为
\begin{choices}[columns=4,label-pos=left]
\item $\dfrac15$．
\item $\dfrac12$．
\item $2$．
\item $3$．
\end{choices}
\end{question}
\QuestionBreak

\begin{question}[index=7]
设 $\symbfit A,\symbfit B$ 为 $n$ 阶矩阵且 $\symbfit A\symbfit B=\symbfit A+\symbfit B$，则以下结论\par
① $(\symbfit A-\symbfit E)^{-1}=\symbfit B-\symbfit E$；\par
② $\symbfit A\symbfit B=\symbfit B\symbfit A$；\par
③ $r(\symbfit A)=r(\symbfit B)$；\par
④ $\symbfit A,\symbfit B$ 的特征向量是公共的．\par
正确结论的个数为
\begin{choices}[columns=4]
\item $4$．
\item $3$．
\item $2$．
\item $1$．
\end{choices}
\end{question}
\QuestionBreak

\begin{question}[index=8]
设 $X_n\ (n\ge1)$ 相互独立且都在 $[0,1]$ 上服从均匀分布，则当 $n\to\infty$ 时，
$\displaystyle \left(\prod_{k=1}^n X_k\right)^{1/n}\xrightarrow{P}$
\begin{choices}[columns=4]
\item $\mathrm e^{-\frac12}$．
\item $\mathrm e^{-1}$．
\item $\mathrm e^{\frac12}$．
\item $\mathrm e$．
\end{choices}
\end{question}
\QuestionBreak

\begin{question}[index=9]
设随机变量 $X\sim N(1,1)$，令 $f(t)=E[(X+t)^2]$，
则 $f(t)$ 在 $[-2,0]$ 上的最小值与最大值分别为
\begin{choices}[columns=4]
\item $1,2$．
\item $1,4$．
\item $-1,2$．
\item $-1,4$．
\end{choices}
\end{question}
\QuestionBreak

\begin{question}[index=10]
设二维随机变量 $\displaystyle (X,Y)\sim N\left(1,1;2,2;\frac12\right)$，则
\begin{choices}[columns=2]
\item $X\sim N(1,1)$．
\item $X+Y\sim N(2,8)$．
\item $\operatorname{Cov}(X,Y)=1$．
\item $E[(X-Y)^2]=6$．
\end{choices}
\end{question}
